% Part: normal-modal-logic
% Chapter: axioms-systems
% Section: systems-distinct

\documentclass[../../../include/open-logic-section]{subfiles}

\begin{document}

\olfileid{nml}{prf}{dis}

\olsection{તંત્રો ભિન્ન છે તે બતાવવું}

\olref[prs]{sec}માં આપણે બે મોડલ તર્ક-તંત્રો હકીકતમાં એક જ તંત્ર છે તે
કેવી રીતે સાબિત કરવું જોયું. \olref[snd]{thm:soundness} વડે બે મોડલ
તંત્રો $\Sigma$ અને $\Sigma'$ ભિન્ન છે તે બતાવી શકાય: એવું
!!a{formula}~$!A$ શોધો કે $\Sigma' \Proves !A$, પરંતુ તે~$\Sigma$ના
કોઈ નિદર્શનમાં અસત્ય હોય.

\begin{prop}
  $\Log{KD} \subsetneq \Log{KT}$
\end{prop}

\begin{proof} આ બધા સ્વવાચક સંબંધો સિરિયલ હોય છે એ અર્થાનુસારી હકીકતનો
  વાક્યરચનાલક્ષી સમકક્ષ છે. $\Log{KD} \subseteq
  \Log{KT}$ બતાવવા $\Log{KD} \Proves !B$ હોય તો
  $\Log{KT} \Proves !B$ થાય છે તે જોવું પડે. તે $\Log{KT} \Proves
  \Ax{D}$ પરથી મળે છે, જેમ
  \olref[prs]{prop:S5facts}\olref[prs]{prop:S5facts-KT-D}માં બતાવ્યું છે.\sourcecorrection{OLMD-028}{મૂળમાં અહીં KTમાંથી 'Log D' નિષ્પન્ન થાય એમ છે; નિષ્પન્ન થતું તંત્ર નહિ પરંતુ Dનું સ્વયંસિદ્ધિ-સૂત્ર છે, તેથી Ax D મૂક્યું છે.}
  સમાવેશ ચુસ્ત છે તે બતાવવા યથાર્થતા
  (\olref[snd]{thm:soundness}) મુજબ \Log{KD}નું એવું નિદર્શન આપવું પૂરતું
  છે જેમાં \Ax{T}, એટલે કે $\Box p \lif p$, અસત્ય હોય (આ સરળ કામ
  અભ્યાસપ્રશ્ન તરીકે છોડ્યું છે). ત્યારે યથાર્થતા મુજબ $\Log{KD}
  \Proves/ \Box p \lif p$.
\end{proof}

\begin{prop}
  $\Log{KB} \neq \Log{K4}$. 
\end{prop}

\begin{proof}
  આપણે એવું સંમિત નિદર્શન રચીએ છીએ જેમાં \Ax{4}નો કોઈ દૃષ્ટાંત અસત્ય
  હોય. તે દૃષ્ટાંત \Log{K4}માં સ્પષ્ટપણે !!{derivable} છે, પરંતુ
  \Log{KB}માં નથી; તેથી $\Log{K4} \nsubseteq \Log{KB}$ ફલિત થશે.
  \olref{fig:Bnot4}નું સંમિત નિદર્શન~$\mModel{M}$ લો. નિદર્શન સંમિત
  હોવાથી~$\mModel{M}$માં \Ax{K} અને~\Ax{B} સત્ય છે
  (અનુક્રમે \olref[syn][sch]{prop:Kvalid} અને
  \olref[frd][acc]{thm:soundschemas} મુજબ). પરંતુ
  $\mSat/{M}{\Box p \lif \Box\Box p}[w_1]$.
\end{proof}

\begin{figure}[htpb]
  \centering
  \begin{tikzpicture}[modal]
    \node[world] (w1) [label=above:\mFalse{p},
    label={below:$\mSat{{}}{\Box p}$\\ $\mSat/{{}}{\Box\Box p}$}] {$w_1$}; 
    \node[world] (w2) [label=above:\mTrue{p},
      label=below:$\mSat/{{}}{\Box p}$, right=of w1]{$w_2$}; 
    \draw[->, bend left] (w1) to (w2); 
    \draw[->, bend left] (w2) to (w1); 
  \end{tikzpicture}
  \caption{\Ax{4}ના દૃષ્ટાંતને અસત્ય બનાવતું સંમિત નિદર્શન.}
  \ollabel{fig:Bnot4}
\end{figure}
 
\begin{thm}\ollabel{thm:KTBnot45}
  $\Log{KTB} \Proves/ \Ax{4}$ અને $\Log{KTB} \Proves/ \Ax{5}$.
\end{thm}\sourcecorrection{OLMD-029}{મૂળ પ્રમેયમાં KTBમાંથી 'Log 4' અને 'Log 5' નિષ્પન્ન નથી એમ લખાયું છે; નિષ્પન્નતાના જમણે તંત્રના નામને બદલે અનુક્રમે 4 અને 5નાં સ્વયંસિદ્ધિ-સૂત્રો મૂક્યાં છે.}

\begin{proof}
  \olref[frd][acc]{thm:soundschemas} મુજબ \Ax{T} અને \Ax{B}ના બધા
  દૃષ્ટાંતો અનુક્રમે દરેક સ્વવાચક અને સંમિત નિદર્શનમાં સત્ય છે. તેથી
  યથાર્થતા મુજબ એવું સ્વવાચક-સંમિત નિદર્શન શોધવું પૂરતું છે જેના કોઈ
  જગતમાં~\Ax{4}નો દૃષ્ટાંત અસત્ય હોય, અને~\Ax{5} માટે પણ એમ જ. બંને
  દાવા માટે આપણે એક જ નિદર્શન વાપરીએ છીએ.
  \olref{fig:KTBnot45}નું સંમિત, સ્વવાચક નિદર્શન લો. ત્યારે
  $\mSat/{M}{\Box p \to \Box \Box
    p}[w_1]$, તેથી \Ax{4} એ $w_1$માં અસત્ય છે. એ જ રીતે,
  $\mSat/{M}{\Diamond
    \lnot p \lif \Box \Diamond \lnot p}[w_2]$, તેથી $!A = \lnot p$
  મૂકેલો~\Ax{5}નો દૃષ્ટાંત~$w_2$માં અસત્ય છે.
\end{proof}

\begin{figure}[htpb]
  \centering
  \begin{tikzpicture}[modal]
    \node[world] (w1) [label={right:\mTrue{p}},
      label={below: $\mSat{{}}{\Box p}$\\
        $\mSat/{{}}{\Box\Box p}$\\
        $\mSat/{{}}{\Diamond\lnot p}$}] {$w_1$};
    \draw[reflexive above] (w1) to (w1);
    \node[world] (w2) [label={right:\mTrue{p}}, label={below:
        $\mSat{{}}{\Diamond\lnot p}$\\
        $\mSat/{{}}{\Box\Diamond\lnot p}$}, right=of w1] {$w_2$} ;
    \draw[reflexive above] (w2) to (w2);
    \node[world] (w3) [label={right:\mFalse{p}},right=of w2] {$w_3$};
    \draw[reflexive above] (w3) to (w3);
    \draw[->, bend left] (w1) to (w2);  
    \draw[->, bend left] (w2) to (w3);  
    \draw[->, bend left] (w3) to (w2);  
    \draw[->, bend left] (w2) to (w1);  
  \end{tikzpicture}
  \caption{\olref{thm:KTBnot45} માટેનું નિદર્શન.}
  \ollabel{fig:KTBnot45}
\end{figure}

\begin{thm}\ollabel{thm:KD5not4}
  $\Log{KD5} \neq \Log{KT4} = \Log{S4}$.
\end{thm}

\begin{proof}
  \olref[frd][acc]{thm:soundschemas} મુજબ \Ax{D} અને~\Ax{5}ના બધા
  દૃષ્ટાંતો દરેક સિરિયલ યુક્લિડિયન નિદર્શનમાં સત્ય છે. તેથી એવું
  સિરિયલ યુક્લિડિયન નિદર્શન શોધવું પૂરતું છે જેના કોઈ જગતમાં~\Ax{4}નો
  દૃષ્ટાંત અસત્ય હોય. \olref{fig:KD5not4}નું નિદર્શન લો અને નોંધો કે
  $\mSat/{M}{\Box p
  \lif \Box\Box p}[w_1]$.
\end{proof}

\begin{figure}[t]
  \centering
  \begin{tikzpicture}[modal]
    \node[world] (w2) [label=north west:\mTrue{p}]{$w_2$} ;
    \draw[reflexive left] (w2) to (w2);
    \node[world] (w1) [label=right:\mFalse{p},
      label=below:{$\mSat{{}}{\Box p}, \mSat/{{}}{\Box\Box p}$},
      below right=of w2]{$w_1$};
    \node[world] (w3) [label=north east:\mTrue{p}, above right=of w1]{$w_3$} ;
    \draw[reflexive right] (w3) to (w3);
    \node[world] (w4) [label=right:\mFalse{p},above right=of w2]{$w_4$};
    \draw[reflexive above] (w4) to (w4);
    \draw[->] (w1) to (w2);  
    \draw[->] (w1) to (w3);  
    \draw[->, bend left=15] (w2) to (w3);
    \draw[->, bend left=15] (w3) to (w2);  
    \draw[->, bend left=15] (w2) to (w4);
    \draw[->, bend left=15] (w4) to (w2);  
    \draw[->, bend left=15] (w3) to (w4);
    \draw[->, bend left=15] (w4) to (w3);  
  \end{tikzpicture}
  \caption{\olref{thm:KD5not4} માટેનું નિદર્શન.}
  \ollabel{fig:KD5not4}
\end{figure}

\begin{prob}
  $3$ જગતોવાળું નિદર્શન વાપરીને \olref[nml][prf][dis]{thm:KD5not4}નું
  વૈકલ્પિક પ્રમાણ આપો.
\end{prob}

\begin{prob}
  $\Log{KT4} \Proves/ \Ax{B}$ અને $\Log{KT4} \Proves/ \Ax{5}$ બંને
  બતાવતું એક જ સ્વવાચક પરંપરિત નિદર્શન આપો.
\end{prob}

\end{document}
